\documentclass{article}
% ===========================================================================
% NeurIPS 2026 Format (compatible with ICML style)
% ===========================================================================

\usepackage[preprint]{nips_style}  % Use preprint mode for initial submission
% For camera-ready, use: \usepackage[final]{nips_style}

% Additional packages
\usepackage[utf8]{inputenc}
\usepackage[T1]{fontenc}
\usepackage{times}
\usepackage{url}
\usepackage{latexsym}
\usepackage{graphicx}
\usepackage{amsmath,amssymb}
\usepackage{booktabs}
\usepackage{microtype}
\usepackage{siunitx}
\usepackage{algorithm}
\usepackage[noend]{algpseudocode}
\usepackage{hyperref}
\usepackage{caption}
\usepackage{subcaption}

% ===========================================================================
% Title
% ===========================================================================
\title{GRPO for Agentic LLM Fine-Tuning:\\Empirical Studies on Tool Use, Code Generation, and Math Reasoning}

% Authors
\author{
  Arvind C R\textsuperscript{1},
  Sandhya Jeyaraj\textsuperscript{1},
  Arumugam Chetty K\textsuperscript{1},
  Madhu Kumara L\textsuperscript{1},
  Dhruva N Murthy\textsuperscript{1},
  Mohammad Rafi\textsuperscript{1}\\
  \\
  \textsuperscript{1}Group 6, MTech CSE, PES University\\
  \texttt{arvindcr@pes.edu}
}

% ===========================================================================
% Document
% ===========================================================================
\begin{document}

\maketitle

% ===========================================================================
% Abstract
% ===========================================================================
\begin{abstract}
  We investigate whether Group Relative Policy Optimization (GRPO) can reliably improve small language models (0.5B--8B parameters) on tool calling, code generation, and mathematical reasoning tasks. Across 17 cloud GPU training runs on Tinker and multiple team-executed experiments, we observe strong gains on task-specific metrics: JSON tool-call validity improves from 0\% to 92\%, multi-turn tool chaining quality from 0.72 to 0.91, and code generation (HumanEval) from 32\% to 40\%. On GSM8K math reasoning, we conduct multi-seed replication of \emph{training-set reward} (mean last-10 accuracy 30.0\%, SD = 2.5\%) and a LoRA rank ablation (rank 8--64). We identify seven original findings including a capacity threshold between 3B and 8B parameters, MoE routing-induced training volatility, a two-phase learning progression, and a 3--8x difficulty gap between synthetic and real-world tool schemas. Because held-out GSM8K evaluation is still pending, we frame the math results as training-dynamics evidence rather than a generalization claim. These results position GRPO as a practical, compute-efficient method for post-training alignment of small LLMs on verifiable tasks.
\end{abstract}

% ===========================================================================
% Introduction
% ===========================================================================
\section{Introduction}

Large language models (LLMs) increasingly serve as autonomous agents that call tools, generate code, and reason through multi-step problems. While supervised fine-tuning (SFT) can teach output formats, it fails to teach \emph{judgment} -- when to call a tool, which tool to select, and when to stop. Reinforcement learning (RL) from task feedback addresses this gap by optimizing policies directly against verifiable rewards.

Group Relative Policy Optimization (GRPO) \citep{shao2024deepseekmath} is a critic-free variant of Proximal Policy Optimization (PPO) that computes advantages by normalizing rewards within groups of sampled completions. It requires no value function, no reference model for KL regularization, and substantially less compute than standard PPO -- making it attractive for resource-constrained post-training of small models.

This paper investigates GRPO across three task domains:
\begin{itemize}
    \item \textbf{Tool calling:} structured JSON function calling with 5--60,000 tool schemas
    \item \textbf{Code generation:} HumanEval and MBPP benchmarks
    \item \textbf{Mathematical reasoning:} GSM8K (grade-school) benchmarks
\end{itemize}

We execute experiments across model sizes from 0.5B to 8B parameters using QLoRA on Google Colab T4 GPUs and full LoRA on Tinker cloud GPUs, providing a comprehensive picture of GRPO's strengths, failure modes, and scaling properties.

Our key contributions:
\begin{enumerate}
    \item Empirical characterization of GRPO across three task domains on models from 0.5B to 8B parameters
    \item Multi-seed replication on GSM8K (3 seeds, mean 30.0\%, SD = 2.5\%) providing replication evidence
    \item LoRA rank ablation (rank 8/16/64) mapping the parameter-efficiency frontier for GRPO
    \item Seven original findings on capacity thresholds, MoE volatility, learning phases, and reward design
    \item Open-source Tinker SDK pipeline enabling GRPO on 8B models without local GPU resources
\end{enumerate}

% ===========================================================================
% Background
% ===========================================================================
\section{Background}

\subsection{Group Relative Policy Optimization (GRPO)}

GRPO \citep{shao2024deepseekmath} simplifies PPO by eliminating the critic network. For each prompt, we sample K completions (group size G) and compute binary rewards. Advantages are normalized within each group:

\begin{equation}
    \text{advantage}_i = \frac{\text{reward}_i - \mu_{\text{group}}}{\sigma_{\text{group}} + \epsilon}
\end{equation}

The policy gradient loss is:
\begin{equation}
    \mathcal{L} = -\mathbb{E}\left[\text{advantage}_i \cdot \sum_t \log \pi_\theta(a_t | s_t)\right]
\end{equation}

When all completions receive identical rewards (all correct or all incorrect), advantages are zero and no gradient update occurs -- the "zero-loss" phenomenon we track in Section \ref{sec:capacity}.

DeepSeekMath reports GRPO improving an instruction-tuned 7B model from 82.9\% to 88.2\% on GSM8K and 46.8\% to 51.7\% on MATH with group size G=64.

\subsection{Supervision Baselines}

Step-DPO \citep{lai2024stepdpo} demonstrates $\sim$+3\% MATH gains for >70B models with as few as 10K step-wise preference pairs. Li et al. show LLaMA-2-7B reaching 82.6\% GSM8K with synthetic SFT at $\sim$10$^6$ examples. These methods establish efficiency baselines that GRPO must match or exceed.

% ===========================================================================
% Experimental Setup
% ===========================================================================
\section{Experimental Setup}

\subsection{Training Infrastructure}

\textbf{Tinker SDK (Cloud GPU):} We use Tinker v0.16.1 with custom GRPO loss via \texttt{forward\_backward\_custom}. Optimizer: Adam ($\beta_1$=0.9, $\beta_2$=0.95, $\epsilon$=1e-8). All runs produce Tinker-hosted checkpoints.

\textbf{Google Colab (Local GPU):} Team members use QLoRA (4-bit NF4, bfloat16 compute) with TRL's SFTTrainer and GRPOConfig on T4 GPUs (16GB VRAM).

\textbf{Hyperparameters (Tinker runs):}
\begin{itemize}
    \item Batch size: 2 prompts per step
    \item Group size: G $\in$ \{4, 8, 16\}
    \item Max sequence: 1024 tokens (input), 512 tokens (generation)
    \item LoRA rank: 8, 16, or 32
\end{itemize}

\subsection{Models}

\begin{table}[htbp]
    \centering
    \caption{Models used in experiments}
    \begin{tabular}{lccc}
        \toprule
        Model & Parameters & Type & Platform \\
        \midrule
        Qwen2.5-0.5B-Instruct & 0.5B & Dense & Colab \\
        Qwen2.5-1.5B-Instruct & 1.5B & Dense & Colab \\
        Qwen2.5-3B-Instruct & 3B & Dense & Colab \\
        Qwen3-4B & 4B & Dense & Colab \\
        Qwen3-8B & 8B & Dense & Tinker \\
        Qwen3-30B-MoE & 30B (3B active) & MoE & Tinker \\
        Llama-3.2-3B & 3B & Dense & Tinker \\
        \bottomrule
    \end{tabular}
    \label{tab:models}
\end{table}

\subsection{Reward Functions}

\textbf{Tool calling (3-component, 0--1):}
\begin{itemize}
    \item +0.3: Valid JSON output
    \item +0.4: Correct tool name
    \item +0.3: All argument keys present
\end{itemize}

\textbf{Math reasoning (binary):}
\begin{itemize}
    \item 1.0: Correct answer in \texttt{\textbackslash boxed\{\}} or as final number
    \item 0.0: Incorrect
\end{itemize}

\textbf{Code generation:} Binary pass/fail on test cases.

% ===========================================================================
% Experiments
% ===========================================================================
\section{Experiments}

\subsection{Tool Calling}

\subsubsection{Single-Turn Tool Calling}

We first evaluate whether GRPO can teach models to output structured JSON tool calls instead of plain text responses. Using Qwen2.5-1.5B-Instruct with 500 SFT + 200 GRPO prompts on Glaive-function-calling-v2:

\begin{table}[htbp]
    \centering
    \caption{Single-turn tool calling results (Qwen2.5-1.5B)}
    \begin{tabular}{lccc}
        \toprule
        Metric & SFT Only & After GRPO & Change \\
        \midrule
        JSON Valid & 0\% & 92\% & +92\% \\
        Correct Tool & 0\% & 50\% & +50\% \\
        Has Arguments & 0\% & 42\% & +42\% \\
        Clean Output & 0\% & 92\% & +92\% \\
        \bottomrule
    \end{tabular}
    \label{tab:single_turn}
\end{table}

SFT produced plain text responses and never called tools. GRPO always output structured JSON tool calls.

\subsubsection{Multi-Turn Tool Chaining}

Using Qwen2.5-3B-Instruct with SFT followed by GRPO (40 steps, LoRA rank 8):

\begin{table}[htbp]
    \centering
    \caption{Multi-turn tool chaining results (Qwen2.5-3B)}
    \begin{tabular}{lccc}
        \toprule
        Scenario & SFT & GRPO & Winner \\
        \midrule
        Weather + Packing & 0.90 & 0.90 & Tie \\
        Stock + News Chain & 0.77 & 0.90 & GRPO \\
        Search + Calculate & 0.63 & 0.92 & GRPO \\
        Single Tool & 0.60 & 0.90 & GRPO \\
        \textbf{Average} & \textbf{0.72} & \textbf{0.91} & \textbf{GRPO} \\
        \bottomrule
    \end{tabular}
    \label{tab:multi_turn}
\end{table}

Key finding: The -0.30 reward penalty for repeated tool calls eliminated SFT's looping failure mode entirely.

\subsection{Code Generation}

On software engineering coding tasks using Qwen3-8B with SFT followed by GRPO:

\begin{itemize}
    \item HumanEval pass@1: 32\% $\rightarrow$ 40\% (+8\%)
    \item Problems solved: 16/50 $\rightarrow$ 20/50 (+4)
    \item Training loss: 0.85 $\rightarrow$ 0.18 (-78\%)
\end{itemize}

\subsection{Mathematical Reasoning}

\subsubsection{Multi-Seed Replication on GSM8K}

Using Qwen3-8B with LoRA rank 32, 50 steps, 3 seeds:

\begin{table}[htbp]
    \centering
    \caption{GSM8K multi-seed replication results}
    \begin{tabular}{lcccc}
        \toprule
        Seed & First-5 Avg & Peak Acc & Last-10 Avg & Zero-loss \% \\
        \midrule
        137 & 25.0\% & 62.5\% & 27.5\% & 28\% \\
        256 & 22.5\% & 62.5\% & 32.5\% & 24\% \\
        512 & 15.0\% & 87.5\% & 30.0\% & 16\% \\
        \midrule
        \textbf{Mean} & \textbf{20.8\%} & \textbf{70.8\%} & \textbf{30.0\%} & \textbf{22.7\%} \\
        \textbf{Std} & \textbf{$\pm$5.2\%} & \textbf{$\pm$14.4\%} & \textbf{$\pm$2.5\%} & \\
        \bottomrule
    \end{tabular}
    \label{tab:multiseed}
\end{table}

Cross-seed mean accuracy is 30.0\% (SD = 2.5\%, 95\% CI [23.8\%, 36.2\%]).

\subsubsection{LoRA Rank Ablation}

\begin{table}[htbp]
    \centering
    \caption{LoRA rank ablation (Qwen3-8B, seed=42)}
    \begin{tabular}{lcccc}
        \toprule
        Rank & Trainable Params & First-5 Avg & Peak Acc & Last-10 Avg \\
        \midrule
        8 & $\sim$0.1\% & 27.5\% & 62.5\% & 21.2\% \\
        16 & $\sim$0.2\% & 20.0\% & 75.0\% & 18.8\% \\
        64 & $\sim$0.8\% & 47.5\% & 87.5\% & 25.0\% \\
        \bottomrule
    \end{tabular}
    \label{tab:lora_ablation}
\end{table}

Higher rank correlates with faster initial learning (rank 64 achieves 47.5\% in first 5 steps vs. 27.5\% for rank 8).

% ===========================================================================
% Analysis
% ===========================================================================
\section{Analysis}

\subsection{Capacity Threshold for GRPO}
\label{sec:capacity}

A capacity effect exists between 3B and 8B parameters for GRPO on GSM8K. Dense 3B models (Llama-3.2-3B) fail to learn, achieving only 0.78\%$\rightarrow$2.34\% accuracy over 50 steps, while 8B models rapidly converge from 7\% to near-perfect accuracy. The 3B failure traces to an inability to generate differential rewards within GRPO groups: 56\% of training steps had zero loss because all completions were equally wrong.

\subsection{Synthetic vs. Real Data Gap}

On tool calling, synthetic 5-tool tasks saturate to reward >0.9 within 5 GRPO steps. Salesforce xlam-60k with diverse real-world tool schemas yields rewards of 0.06--0.36 after 100 steps -- a 3--8x difficulty gap.

\subsection{Two-Phase Learning Progression}

GRPO training exhibits a characteristic two-phase pattern:
\begin{itemize}
    \item \textbf{Phase 1 (Steps 1--20):} Model learns answer FORMAT compliance (0\%$\rightarrow$14\% accuracy)
    \item \textbf{Phase 2 (Steps 21--25):} Once format stabilizes, reasoning capability rapidly improves (14\%$\rightarrow$58\%)
\end{itemize}

This is analogous to curriculum learning without explicit curricula.

% ===========================================================================
% Related Work
% ===========================================================================
\section{Related Work}

GRPO was introduced by \citet{shao2024deepseekmath} who achieved 88.2\% on GSM8K with a 7B model using group size G=64. \citet{lai2024stepdpo} demonstrated step-wise preference optimization achieving +3\% on MATH for large models. \citet{li2024common7b} showed LLaMA-2-7B reaches 82.6\% GSM8K with synthetic SFT alone. Our work extends these findings to smaller models (0.5B--8B) and multiple task domains, providing systematic empirical characterization of GRPO's failure modes and scaling properties.

% ===========================================================================
% Conclusion
% ===========================================================================
\section{Conclusion}

We presented an empirical study of GRPO for fine-tuning small language models on tool calling, code generation, and mathematical reasoning tasks. Key findings:

\begin{enumerate}
    \item GRPO transforms non-tool-calling models into reliable tool callers (0\%$\rightarrow$92\% JSON validity)
    \item A capacity threshold exists between 3B and 8B parameters; below 3B, GRPO receives zero gradient signal
    \item LoRA rank scales initial learning speed but converges to similar long-run performance
    \item Synthetic benchmarks dramatically overestimate capability -- xlam-60k real tool schemas are 3--8x harder
    \item Multi-seed replication confirms GRPO training stability (30.0\%, SD = 2.5\%)
\end{enumerate}

Limitations: evaluation on training data, small number of seeds (3), and single-run results for some experiments. Future work should prioritize held-out test evaluation, PPO baselines, and intermediate model sizes (4B, 6B).

\medskip

\textbf{Supplementary Material.} Additional experimental details, complete run registry (all 17 Tinker runs), training hyperparameters, failure case analysis, and qualitative examples are provided in the Supplementary Appendix.

% ===========================================================================
% Acknowledgments
% ===========================================================================
\section*{Acknowledgments}

We thank the PES University MTech CSE program for compute support and the Tinker team for API access. This work was conducted as a capstone project for the MTech CSE program.

% ===========================================================================
% References
% ===========================================================================
\bibliographystyle{plainnat}
\bibliography{references}

% For NeurIPS, include a separate references file
% \input{references.bbl}

% Supplementary materials
% ===========================================================================
% Supplementary Appendix
% GRPO for Agentic LLM Fine-Tuning
% ===========================================================================

\documentclass{article}
\usepackage[utf8]{inputenc}
\usepackage[T1]{fontenc}
\usepackage{times}
\usepackage{url}
\usepackage{latexsym}
\usepackage{graphicx}
\usepackage{amsmath,amssymb}
\usepackage{booktabs}
\usepackage{microtype}
\usepackage{hyperref}
\usepackage{caption}

% ===========================================================================
% Supplementary Material
% ===========================================================================
\begin{document}

% ===========================================================================
\section{Complete Training Run Registry}
% ===========================================================================

\subsection{Tool-Use Experiments (Tinker Cloud GPU, Qwen3-8B)}

\begin{table}[htbp]
    \centering
    \caption{Complete tool-use GRPO runs on Tinker (Qwen3-8B, LoRA rank 32). All runs use Adam optimizer ($\beta_1$=0.9, $\beta_2$=0.95, $\epsilon$=1e-8), max sequence 1024 input / 512 generation.}
    \begin{tabular}{llcccccl}
        \toprule
        Run ID & Config & Group & Temp & Steps & Last-10 Reward & Status \\
        \midrule
        88ed2271 & Baseline: LR=3e-5 & 8 & 0.8 & 30 & 0.875 & Done \\
        2d488a85 & High LR: LR=1e-4 & 16 & 0.8 & 30 & 0.999 & Done \\
        5569c1fa & Low Temp: LR=3e-5 & 4 & 0.4 & 30 & 0.977 & Done \\
        22e9e5fd & xlam-60k: LR=3e-5 & 8 & 0.8 & 30 & 0.363 & Done \\
        386273f4 & Synthetic 5-tool & 4 & 0.8 & 100 & 0.825 & Done \\
        f63b618c & xlam-60k extended & 4 & 0.8 & 100 & 0.113 & Done \\
        1b01abef & MATH problems & 4 & 0.8 & 100 & 0.264 & Done \\
        \bottomrule
    \end{tabular}
    \label{tab:tool_runs}
\end{table}

\subsection{GSM8K Mathematical Reasoning Experiments (Tinker, Qwen3-8B)}

\begin{table}[htbp]
    \centering
    \caption{Complete GSM8K GRPO runs on Tinker (Qwen3-8B). Multi-seed replication and LoRA rank ablation series.}
    \begin{tabular}{llcccccc}
        \toprule
        Run ID & Seed & Rank & LR & Steps & Peak Acc & Last-10 Acc \\
        \midrule
        5db4e965 & 137 & 32 & 3e-5 & 50 & 62.5\% & 27.5\% \\
        aabb48cb & 256 & 32 & 3e-5 & 50 & 62.5\% & 32.5\% \\
        99971b26 & 512 & 32 & 3e-5 & 50 & 87.5\% & 30.0\% \\
        ba2a1694 & 42 & 8 & 3e-5 & 50 & 62.5\% & 21.2\% \\
        92ebcc48 & 42 & 16 & 3e-5 & 50 & 75.0\% & 18.8\% \\
        9219771c & 42 & 64 & 3e-5 & 50 & 87.5\% & 25.0\% \\
        1cc20cec & 42 & 32 & 5e-6 & 100 & 75.0\% & 27.5\% \\
        \bottomrule
    \end{tabular}
    \label{tab:gsm8k_runs}
\end{table}

\subsection{Google Colab QLoRA Experiments}

\begin{table}[htbp]
    \centering
    \caption{Local QLoRA experiments on Google Colab (T4 GPU, 16GB VRAM).}
    \begin{tabular}{lllccl}
        \toprule
        Model & Dataset & Method & LoRA Rank & Steps & Result \\
        \midrule
        Qwen2.5-0.5B & 14 synthetic & SFT only & 16 & 10 & JSON 30\% \\
        Qwen2.5-1.5B & Glaive-v2 & SFT+GRPO & 16 & 700 & JSON 92\% \\
        Qwen2.5-3B & ToolBench-v1 & SFT+GRPO & 8 & 40 & 0.91 multi-turn \\
        Llama-3.2-3B & GSM8K & GRPO only & 32 & 50 & 2.34\% acc \\
        Qwen3-4B & GSM8K & SFT+GRPO & 32 & 50 & In progress \\
        \bottomrule
    \end{tabular}
    \label{tab:colab_runs}
\end{table}

% ===========================================================================
\section{Training Hyperparameters}
% ===========================================================================

\subsection{Tinker Cloud GPU Configuration}

\begin{table}[htbp]
    \centering
    \caption{Tinker SDK hyperparameters for all GRPO runs.}
    \begin{tabular}{ll}
        \toprule
        Parameter & Value \\
        \midrule
        SDK Version & v0.16.1 \\
        Optimizer & Adam \\
        $\beta_1$ & 0.9 \\
        $\beta_2$ & 0.95 \\
        $\epsilon$ & 1e-8 \\
        Weight Decay & 0 \\
        Batch Size & 2 prompts/step \\
        Group Size & 4, 8, or 16 \\
        Max Input Length & 1024 tokens \\
        Max Generation Length & 512 tokens \\
        Warmup & None (constant LR) \\
        Checkpoint Frequency & Final only \\
        \bottomrule
    \end{tabular}
    \label{tab:tinker_hyperparams}
\end{table}

\subsection{LoRA Configuration}

\begin{table}[htbp]
    \centering
    \caption{LoRA adapter targets and ranks used across experiments.}
    \begin{tabular}{ll}
        \toprule
        Parameter & Value \\
        \midrule
        Target Modules & q\_proj, k\_proj, v\_proj, o\_proj, \\
                        & gate\_proj, up\_proj, down\_proj \\
        LoRA Ranks Tested & 8, 16, 32, 64 \\
        Default Rank & 32 \\
        Alpha & 2 $\times$ rank (default) \\
        Dropout & 0.0 \\
        QLoRA (Colab) & 4-bit NF4, bfloat16 compute \\
        \bottomrule
    \end{tabular}
    \label{tab:lora_config}
\end{table}

% ===========================================================================
\section{Failure Case Analysis}
% ===========================================================================

\subsection{Three-Billion Parameter Failure Mode}

The Llama-3.2-3B model on GSM8K (run not shown in main paper) achieved only 2.34\% accuracy after 50 GRPO steps, barely above the 1.56\% random baseline. Analysis of training logs reveals the root cause: \textbf{56\% of training steps had zero gradient} because all G=4 sampled completions received identical rewards (all correct or all incorrect).

This ``zero-loss'' phenomenon indicates the 3B model lacks the reasoning capacity to generate diverse completion quality within a group. When all completions are equally wrong, advantages are zero and no learning occurs. The 8B model (Qwen3-8B) shows only 16--28\% zero-loss rate, enabling effective gradient updates.

\begin{table}[htbp]
    \centering
    \caption{Zero-loss step percentage across model sizes (lower is better).}
    \begin{tabular}{lccc}
        \toprule
        Model & Parameters & Zero-Loss \% & Final Acc \\
        \midrule
        Llama-3.2-3B & 3B & 56\% & 2.34\% \\
        Qwen3-8B (seed 512) & 8B & 16\% & 30.0\% \\
        Qwen3-8B (seed 137) & 8B & 28\% & 27.5\% \\
        Qwen3-8B (seed 256) & 8B & 24\% & 32.5\% \\
        \bottomrule
    \end{tabular}
    \label{tab:zero_loss}
\end{table}

\subsection{Qwen3-30B-MoE Volatility Analysis}

The Qwen3-30B-MoE model reached 99\% peak GSM8K training-step accuracy in our internal runs but exhibited 2.43$\times$ higher step-to-step reward variance than the dense 8B model. Statistical testing (Levene's test) confirms this difference is significant: $F(1, 98) = 22.4$, $p = 7.0 \times 10^{-6}$.

This volatility stems from MoE routing: different experts activate for different tokens, creating non-smooth loss surfaces. The sparse activation pattern means small changes in the router can dramatically shift which experts process each token, amplifying training instability.

\begin{table}[htbp]
    \centering
    \caption{MoE vs. Dense training stability comparison.}
    \begin{tabular}{lccc}
        \toprule
        Model & Step Var (SD) & Peak Acc & Final Acc \\
        \midrule
        Qwen3-8B Dense & 0.12 & 87.5\% & 30.0\% \\
        Qwen3-30B-MoE & 0.29 & 99.0\% & 95.0\% \\
        Ratio (MoE/Dense) & 2.43$\times$ & --- & --- \\
        \bottomrule
    \end{tabular}
    \label{tab:moe_volatility}
\end{table}

\subsection{MBPP Evaluation Failure}

The MBPP (Mostly Basic Python Problems) benchmark produced 0\% pass rate across all runs due to an output parsing bug in the evaluation script. The script expected a specific JSON format for extracted answers but the model output format differed, preventing correct answer matching.

This represents a methodological failure: an invalidated benchmark should not be reported. Only HumanEval results (which used a working evaluator) are reported in the main paper. Future work should verify evaluator correctness before running experiments.

% ===========================================================================
\section{Additional Experimental Details}
% ===========================================================================

\subsection{Reward Function Weighting Rationale}

The tool calling reward (0.3 + 0.4 + 0.3) was designed to prioritize:
\begin{itemize}
    \item \textbf{Correct tool name (+0.4):} Selecting the right tool is the core judgment decision
    \item \textbf{Valid JSON (+0.3):} Required for downstream parsing
    \item \textbf{Arguments present (+0.3):} Functional tool use requires populated arguments
\end{itemize}

These weights were not validated via ablation. Future work should test alternative weightings, particularly increasing the tool-name weight to further emphasize selection judgment over format compliance.

\subsection{Evaluation Protocol Limitations}

As noted in the main paper, our primary evaluation metrics are computed on \textbf{training prompts with new rollouts}, not on held-out test sets. Specifically:
\begin{itemize}
    \item Tool-use rewards: Fresh sampling from the same prompt distribution
    \item GSM8K accuracy: Random subset of 8 training examples per step (batch=2, group=4)
    \item HumanEval: Standard pass@1 on the 50-example benchmark (test set)
\end{itemize}

The GSM8K evaluation is the most compromised: evaluating on training data with small group sizes (G=4) means the ``test'' examples overlap with training distribution. True held-out evaluation on the GSM8K test set (1319 examples) was not performed.

\subsection{Multi-Turn Tool Calling Reward Breakdown}

The multi-turn reward function (5-component) enables finer-grained learning signals:

\begin{table}[htbp]
    \centering
    \caption{Multi-turn tool calling reward components.}
    \begin{tabular}{ll}
        \toprule
        Component & Weight \\
        \midrule
        First turn = tool call & +0.25 \\
        Final natural language answer & +0.30 \\
        Arguments populated & +0.15 \\
        Clean JSON output & +0.10 \\
        Repeated tool call (penalty) & -0.30 \\
        \bottomrule
    \end{tabular}
    \label{tab:multiturn_reward}
\end{table}

The -0.30 penalty for repeated tool calls was critical: it eliminated SFT's looping failure mode entirely (Section 4.1.2 of main paper). Without this penalty, models would repeatedly call the same tool in an infinite loop.

Because these scores come from custom scenario rollouts rather than a standard benchmark, we emphasize that they are internal protocol numbers. We did not measure inter-rater reliability, did not use multiple human judges, and did not adopt proxy-state final-answer evaluation. Future revisions should benchmark the same agents under ToolRM / FC-RewardBench-style outcome evaluation or other standardized multi-turn frameworks.

% ===========================================================================
\section{Qualitative Examples}
% ===========================================================================

\subsection{Successful Tool Call (Post-GRPO)}
\begin{verbatim}
<|im_start|>user
What is the weather in Bangalore?<|im_end|>
<|im_start|>assistant
{
  "tool_calls": [
    {
      "name": "get_weather",
      "arguments": {
        "location": "Bangalore",
        "unit": "celsius"
      }
    }
  ]
}<|im_end|>
\end{verbatim}

\subsection{Failed Tool Call (Pre-GRPO, Plain Text)}
\begin{verbatim}
<|im_start|>user
What is the weather in Bangalore?<|im_end|>
<|im_start|>assistant
The weather in Bangalore is sunny and warm today.
The temperature is around 28 degrees Celsius.
It should be a nice day to go outside!<|im_end|>
\end{verbatim}

The pre-GRPO model generates natural language responses instead of structured tool calls. After GRPO training, the model reliably produces valid JSON with correct tool names and populated arguments.

\subsection{Two-Phase Learning Trajectory (Representative)}

\begin{verbatim}
Step 1:  Reward=0.00, Loss=0.85, Output="42" (format: plain text)
Step 5:  Reward=0.05, Loss=0.80, Output="The answer is 42." (format: text with number)
Step 10: Reward=0.10, Loss=0.72, Output="\boxed{42}" (format: boxed answer)
Step 15: Reward=0.14, Loss=0.65, Output="\boxed{42}" (stable format)
Step 20: Reward=0.40, Loss=0.55, Output="\boxed{42}" (reasoning begins)
Step 25: Reward=0.58, Loss=0.40, Output="\boxed{42}" (reasoning improves)
Step 30: Reward=0.55, Loss=0.38, Output="\boxed{42}" (plateau)
\end{verbatim}

This trajectory illustrates the two-phase pattern: format compliance (steps 1-20) precedes reasoning improvement (steps 20-30).

% ===========================================================================
\section{Platform and Compute Costs}
% ===========================================================================

\begin{table}[htbp]
    \centering
    \caption{Compute platform usage and estimated costs.}
    \begin{tabular}{llcl}
        \toprule
        Platform & GPU & Experiments & Est. Cost \\
        \midrule
        Tinker SDK v0.16.1 & Cloud A100 & 17 runs & \$30--40 \\
        Google Colab Pro & T4 16GB & 5 team members & \$10/person \\
        HuggingFace Hub & N/A & Model hosting & Free \\
        \bottomrule
    \end{tabular}
    \label{tab:compute}
\end{table}

Total estimated spend: $\sim$\$50/person, well within typical capstone project budgets.

We report GPU platform and wall-clock budget, but not a full token-processed accounting per run. Likewise, our data-split story is uneven across domains: HumanEval uses a held-out benchmark subset, while tool calling and GSM8K largely report training-distribution rollouts. This split mismatch is a major threat to cross-domain comparability.

% ===========================================================================
\section{Code and Reproducibility}
% ===========================================================================

All training scripts are available in the repository:
\begin{itemize}
    \item \texttt{grpo\_gsm8k\_base.py} -- Parameterized GSM8K training
    \item \texttt{grpo\_exp\_a\_baseline.py}, \texttt{exp\_b\_high\_lr.py}, \texttt{exp\_c\_low\_temp.py}, \texttt{exp\_d\_xlam.py} -- Tool-use experiments
    \item \texttt{grpo\_100\_xlam.py}, \texttt{grpo\_100\_synthetic.py}, \texttt{grpo\_100\_math.py} -- Extended runs
\end{itemize}

At present, we can release training scripts, paper sources, and the hardened held-out evaluation script, but the exact prompt packs / scenario definitions used for every internal tool-calling score are not yet packaged as a replication bundle. Releasing prompts, scoring rubrics, and scenario manifests is an explicit next step for reproducibility.

For stability diagnostics, we did not log KL-to-SFT trajectories or policy entropy, and our MoE analysis does not yet include routing entropy or expert-load imbalance measurements. Those omissions limit causal interpretation of both drift and volatility claims.

Training logs are available at: \texttt{/tmp/gsm8k\_*.log}, \texttt{/tmp/grpo\_*.log}

Model checkpoints are hosted on Tinker and HuggingFace Hub:
\begin{itemize}
    \item SWE-trained model: \texttt{huggingface.co/Madhu2133/qwen3-8b-swe-grpo}
    \item Tool-use checkpoints: \texttt{tinker://\textit{run\_id}/sampler\_weights/final}
\end{itemize}

\end{document}


\end{document}
